\section{Introduction}
\label{sec:introduction}

A trend smoother that reconstructs historical observations well
need not produce an accurate forecast. In finite-difference
penalized least squares (PLS), the smoothing penalty trades
fidelity to the past against regularity of the fitted path.
Once the path is extrapolated, however, its terminal
level and differences determine the forecast. Selecting
the penalty to recover a historical trend and selecting
it for future extrapolation are different problems.

This paper investigates a specific predictive criterion.
For a historical window $x_t$, let $H(S)$ be the PLS
smoothing matrix indexed by normalized smoothness
$S\in[0,1]$, and let $G_{d,h}$ continue the fitted
trend $h$ observations into the future. At the final
forecast origin $T$, choose
\begin{equation}
\label{eq:intro_pooled_loss}
F^{\mathrm{pool}}_{T,d,L,h}(S)
=\frac{1}{Mh}\sum_{m=1}^M
\left\|y_{t_m+1:t_m+h}
-G_{d,h}H(S)x_{t_m}\right\|_2^2,
\end{equation}
where each historical validation block has ended
before the final origin. The smoothness decision is
\begin{equation}
\label{eq:intro_pooled_choice}
\widehat S^{\mathrm{FCV}}_{T,d,L,h}
\in\operatorname*{arg\,min}_{S\in[0,1]}
F^{\mathrm{pool}}_{T,d,L,h}(S).
\end{equation}
The final trend is refitted using the newest $L$
observations and extrapolated into an untouched
future test block. Thus the validation objective
is prospective: no future test observation is
available when the operational trend is chosen.

Normalized smoothness is an important computational
coordinate. It transforms the unbounded penalty
range $\lambda\in[0,\infty]$ into $S\in[0,1]$
without requiring an arbitrary upper cutoff.
At $S=0$ the historical observations are unsmoothed;
at $S=1$ the fit is projected onto the polynomial
null space of the difference penalty. Intermediate
values change the terminal differences of the fitted
trend and consequently the polynomial forecast
coefficients. For fixed difference order $d$, this
does not select the polynomial degree, which
remains at most $d-1$; it selects the smoothness
of the trend that is to be continued.

The ingredients themselves are established.
\citet{guerrero2007penalized,guerrero2008smoothness}
developed PLS and the percentage-of-smoothness
interpretation, while \citet{cortes2017smoothness}
examined attained smoothness under conventional
criteria in this journal. Controlled-smoothness
trends were already used for remittance and
exchange-rate forecasting
\citep{islas2019remittances,islas2025exchange}.
More directly, \citet{hart1994tscv} used
one-step-ahead predictive error to tune kernel
smoothing, and \citet{vilar2007nonparametric}
studied nonparametric time-series forecasting.
We therefore do not claim that predictive
smoothness tuning in general, or the index
itself, originated here. The statistical
object examined is their concrete combination
into horizon-matched \emph{future-block}
selection of normalized PLS trend smoothness.

A Monte Carlo study addresses whether matching
the validation horizon to the operational
forecast horizon changes smoothness and
out-of-sample error; how that choice compares
with ordinary CV, GCV, AICc, and
simulation-only recovery oracles; and how
the bounded optimization objective behaves
when it contains several minima. The frozen
study crosses five trend mechanisms, three
noise families, two noise levels, and 100
seeds, with six outer origins and four
forecast horizons, yielding 72,000 decisions.
These evaluations are repeated within seeds
and are not 72,000 independent replications.

The main result is an aggregate improvement
of horizon-matched tuning over one-step
tuning at longer horizons for this design.
Neither the mathematical reparameterization
nor the experiment establishes uniform
forecast dominance under all time-series
models. 